
\documentclass[11pt,a4paper]{article}

\usepackage[utf8]{inputenc}
\usepackage[T1]{fontenc}
\usepackage[french]{babel}
\usepackage[a4paper,margin=2cm]{geometry}
\usepackage{amsmath,amssymb,amsthm,mathrsfs}
\usepackage{lmodern}
\usepackage{enumitem}
\usepackage{xcolor}
\usepackage{fancyhdr}
\usepackage{hyperref}
\usepackage{tikz}
\usepackage{siunitx}
\usepackage[version=4]{mhchem}
\usepackage{physics}

\hypersetup{
    colorlinks=true,
    linkcolor=black,
    urlcolor=blue
}

\setlength{\parindent}{0pt}
\setlength{\parskip}{0.5em}

\pagestyle{fancy}
\fancyhf{}
\fancyhead[L]{TD-problème --- Physique-Chimie Terminale}
\fancyhead[R]{Niveau bac renforcé / transition supérieur}
\fancyfoot[C]{\thepage}

\newcommand{\etoile}{\ensuremath{\star}}
\newcommand{\e}{\mathrm{e}}
\newcommand{\R}{\mathbb{R}}

\begin{document}

\begin{center}
    {\LARGE \textbf{Fiche de gros problèmes --- Physique-Chimie Terminale}}\\[0.4em]
    {\large Sujets de type bac, plus exigeants, centrés sur la compréhension profonde des notions}
\end{center}

\textbf{Objectif de la fiche.} Cette fiche ne cherche pas à multiplier les applications directes. Elle propose au contraire de \textbf{gros exercices structurés}, en plusieurs parties, dans l'esprit des sujets de bac, mais avec un niveau de réflexion plus élevé : interprétation physique, critique des modèles, discussion des hypothèses, exploitation de documents, liens entre plusieurs chapitres.

\textbf{Niveaux de difficulté.}
\begin{itemize}[leftmargin=1.5em]
    \item [\etoile\etoile\etoile] problème exigeant de niveau bac solide ;
    \item [\etoile\etoile\etoile\etoile] problème difficile demandant une bonne autonomie ;
    \item [\etoile\etoile\etoile\etoile\etoile] problème très difficile, de synthèse, proche d'une transition vers le supérieur.
\end{itemize}

\textbf{Mode d'emploi.}
\begin{itemize}[leftmargin=1.5em]
    \item Chaque problème doit être rédigé proprement, comme en devoir surveillé.
    \item On attend de l'élève qu'il justifie les modèles utilisés, pas seulement qu'il applique des formules.
    \item Certains résultats numériques ne sont pas le but principal : le sens physique et chimique compte autant que le calcul.
\end{itemize}

\hrule

\section*{Problème 1 --- Freinage d'une trottinette électrique en ville \hfill [\etoile\etoile\etoile]}

Une trottinette de masse totale $m=95\ \si{kg}$ (conducteur compris) roule en ligne droite sur une route horizontale à la vitesse de $\SI{25}{km.h^{-1}}$.
Le conducteur freine pour s'arrêter avant un feu rouge.
On suppose dans un premier temps que la force de freinage est constante et opposée au mouvement.

\begin{enumerate}[label=\arabic*)]
    \item Convertir la vitesse initiale en unités SI.
    \item Rappeler la relation fondamentale de la dynamique dans le référentiel terrestre supposé galiléen.
    \item Montrer que l'accélération du système est constante si la force de freinage est constante.
    \item On estime la valeur de la force de freinage à $\SI{190}{N}$. Déterminer l'accélération.
    \item En déduire la durée du freinage puis la distance de freinage.
    \item Expliquer pourquoi la distance d'arrêt réelle du véhicule est supérieure à la distance de freinage calculée.
    \item Dans la réalité, la force de freinage n'est pas strictement constante. Expliquer qualitativement comment cela modifie le modèle précédent.
    \item Discuter les hypothèses du modèle : route horizontale, absence de vent, système assimilé à un point matériel, frottements regroupés dans une seule force.
\end{enumerate}

\section*{Problème 2 --- Chute d'une balle : modèle idéal et modèle réel \hfill [\etoile\etoile\etoile\etoile]}

On lâche sans vitesse initiale une balle de tennis depuis un balcon situé à $\SI{12,0}{m}$ au-dessus du sol.

\textbf{Partie A --- Modèle de chute libre}
\begin{enumerate}[label=\arabic*)]
    \item Définir la chute libre.
    \item Établir l'expression de la vitesse en fonction du temps dans ce modèle.
    \item Établir l'expression de l'altitude en fonction du temps.
    \item Déterminer la durée théorique de la chute.
    \item Déterminer la vitesse juste avant l'impact.
\end{enumerate}

\textbf{Partie B --- Discussion du modèle}
\begin{enumerate}[label=\arabic*)]
    \setcounter{enumi}{5}
    \item Expliquer pourquoi le modèle précédent n'est qu'une approximation.
    \item Quel est le rôle des frottements de l'air ?
    \item Expliquer qualitativement la notion de vitesse limite.
    \item Dans quel cas l'erreur commise en négligeant l'air est-elle faible ? importante ?
\end{enumerate}

\textbf{Partie C --- Comparaison avec une autre balle}
\begin{enumerate}[label=\arabic*)]
    \setcounter{enumi}{9}
    \item On recommence avec une feuille de papier froissée puis non froissée. Expliquer les différences observées.
    \item Montrer que la masse n'est pas le seul paramètre pertinent pour décrire l'effet de l'air.
\end{enumerate}

\section*{Problème 3 --- Saut à ski et conservation de l'énergie \hfill [\etoile\etoile\etoile]}

Un skieur part sans vitesse initiale d'un point A situé à une altitude supérieure de $\SI{18}{m}$ au point B. On néglige dans un premier temps les frottements.

\begin{enumerate}[label=\arabic*)]
    \item Écrire l'expression de l'énergie mécanique du système skieur-Terre.
    \item Expliquer pourquoi cette énergie se conserve dans le modèle proposé.
    \item Déterminer la vitesse du skieur au point B.
    \item La vitesse réelle mesurée est plus faible que celle obtenue. Que peut-on en conclure ?
    \item En supposant que la perte d'énergie mécanique vaut $\SI{2,5e3}{J}$, déterminer la vitesse réelle du skieur de masse $m=\SI{75}{kg}$ au point B.
    \item Expliquer sous quelles formes peut être dissipée l'énergie mécanique perdue.
    \item Pourquoi dire que l'énergie « disparaît » serait physiquement faux ?
\end{enumerate}

\section*{Problème 4 --- Étude critique d'un airbag \hfill [\etoile\etoile\etoile\etoile]}

Lors d'un choc, un passager passe d'une vitesse de $\SI{50}{km.h^{-1}}$ à zéro.
Sans airbag, l'arrêt se produit sur une distance très courte de l'ordre de quelques millimètres ; avec airbag, la distance d'arrêt est beaucoup plus grande.

\begin{enumerate}[label=\arabic*)]
    \item Convertir $\SI{50}{km.h^{-1}}$ en m\,s$^{-1}$.
    \item Expliquer pourquoi augmenter la durée ou la distance d'arrêt diminue la valeur de la force moyenne subie.
    \item En utilisant un modèle à accélération constante, comparer qualitativement les cas « avec airbag » et « sans airbag ».
    \item Pourquoi un raisonnement uniquement énergétique ne suffit-il pas à rendre compte du danger pour le passager ?
    \item Expliquer en quoi ce problème fait intervenir simultanément la quantité de mouvement, l'énergie et la durée du choc.
    \item Discuter les limites d'un modèle où le passager serait assimilé à un point matériel.
\end{enumerate}

\section*{Problème 5 --- Mesurer la célérité du son \hfill [\etoile\etoile\etoile]}

On veut mesurer la célérité du son dans l'air par une méthode expérimentale simple. On dispose de deux microphones séparés d'une distance connue $d=\SI{2,50}{m}$, d'un système d'acquisition, et d'une source sonore brève.

\begin{enumerate}[label=\arabic*)]
    \item Proposer un protocole expérimental.
    \item Expliquer quelle grandeur doit être mesurée.
    \item Établir la relation permettant de calculer la célérité du son.
    \item Si le retard mesuré vaut $\Delta t = \SI{7,3e-3}{s}$, calculer la célérité du son.
    \item Discuter la précision du résultat.
    \item Citer au moins trois sources d'incertitudes expérimentales.
    \item Expliquer pourquoi il faut veiller à l'alignement correct des microphones et de la source.
\end{enumerate}

\section*{Problème 6 --- Diffraction par une fente et validité du modèle ondulatoire \hfill [\etoile\etoile\etoile\etoile]}

Un laser de longueur d'onde $\lambda = \SI{650}{nm}$ éclaire une fente fine de largeur $a$. L'écran d'observation est placé à une distance $D=\SI{2,0}{m}$.

\begin{enumerate}[label=\arabic*)]
    \item Décrire qualitativement la figure observée.
    \item Expliquer pourquoi la diffraction met en évidence le caractère ondulatoire de la lumière.
    \item On admet que la largeur $L$ de la tache centrale vérifie approximativement $L \propto \dfrac{\lambda D}{a}$. Comment évolue $L$ si $a$ diminue ? interpréter physiquement.
    \item Pourquoi la diffraction est-elle peu visible avec des objets de grande taille devant de la lumière visible ?
    \item On remplace le laser par une lumière blanche. Expliquer qualitativement ce qui change.
    \item Discuter l'intérêt d'utiliser une source monochromatique.
\end{enumerate}

\section*{Problème 7 --- Lunette astronomique et limites d'observation \hfill [\etoile\etoile\etoile\etoile]}

Une lunette astronomique est utilisée pour observer une planète. L'instrument permet de grossir l'image, mais l'observation reste limitée.

\begin{enumerate}[label=\arabic*)]
    \item Expliquer la différence entre « grossir » et « mieux voir ». 
    \item Pourquoi augmenter le diamètre de l'objectif améliore-t-il généralement le pouvoir séparateur ?
    \item Quel lien peut-on faire avec la diffraction ?
    \item Pourquoi l'atmosphère terrestre peut-elle limiter la qualité de l'observation ?
    \item Expliquer pourquoi une image plus grande n'est pas forcément une image plus détaillée.
\end{enumerate}

\section*{Problème 8 --- Circuit de charge d'un condensateur \hfill [\etoile\etoile\etoile]}

On réalise un circuit RC série alimenté par un générateur idéal de tension continue $E=\SI{6,0}{V}$. On utilise une résistance $R=\SI{47}{k\Omega}$ et un condensateur de capacité $C=\SI{10}{\micro F}$.

\begin{enumerate}[label=\arabic*)]
    \item Définir la constante de temps du dipôle RC.
    \item Calculer sa valeur numérique.
    \item Expliquer ce que représente physiquement cette grandeur.
    \item Décrire qualitativement l'évolution de la tension aux bornes du condensateur au cours de la charge.
    \item Pourquoi cette tension ne peut-elle pas varier brutalement ?
    \item Au bout de combien de constantes de temps peut-on considérer la charge comme pratiquement terminée ?
    \item Si l'on double la résistance, quel est l'effet sur la durée de charge ?
    \item Si l'on observe le phénomène à l'oscilloscope, quel réglage doit-on adapter pour visualiser correctement toute la charge ?
\end{enumerate}

\section*{Problème 9 --- Choisir un fusible pour protéger une installation \hfill [\etoile\etoile\etoile\etoile]}

On souhaite protéger un circuit alimentant successivement plusieurs appareils domestiques.

\begin{center}
\begin{tabular}{|c|c|}
\hline
Appareil & Puissance nominale \\
\hline
Bouilloire & \SI{1800}{W} \\
Grille-pain & \SI{900}{W} \\
Lampe & \SI{60}{W} \\
Chargeur & \SI{40}{W} \\
\hline
\end{tabular}
\end{center}

La tension du secteur est prise égale à $U=\SI{230}{V}$.

\begin{enumerate}[label=\arabic*)]
    \item Déterminer l'intensité appelée par chaque appareil lorsqu'il fonctionne seul.
    \item Déterminer l'intensité totale si tous les appareils fonctionnent en même temps.
    \item Expliquer le rôle d'un fusible ou d'un disjoncteur.
    \item Pourquoi protéger une installation en se basant sur la seule puissance d'un appareil isolé serait-il insuffisant ?
    \item Discuter le danger d'une multiprise surchargée.
    \item Pourquoi un court-circuit provoque-t-il un courant très intense ?
\end{enumerate}

\section*{Problème 10 --- Dosage d'un vinaigre commercial \hfill [\etoile\etoile\etoile]}

On veut déterminer la concentration en acide éthanoïque d'un vinaigre commercial. On prélève un volume $V_A=\SI{10,0}{mL}$ de vinaigre dilué, que l'on dose par une solution d'hydroxyde de sodium de concentration $C_B=\SI{0,100}{mol.L^{-1}}$. L'équivalence est obtenue pour $V_E=\SI{13,6}{mL}$.

\begin{enumerate}[label=\arabic*)]
    \item Écrire l'équation de la réaction support du dosage.
    \item Définir l'équivalence.
    \item Établir la relation entre les quantités de matière à l'équivalence.
    \item Déterminer la concentration du vinaigre dilué.
    \item Si le vinaigre commercial a été dilué 10 fois avant dosage, déterminer la concentration initiale.
    \item Expliquer pourquoi il est souvent préférable de diluer le vinaigre avant le dosage.
    \item Décrire l'allure générale de la courbe de dosage pH-métrique.
    \item Pourquoi un dosage d'un acide faible par une base forte ne donne-t-il pas exactement la même courbe qu'un dosage acide fort/base forte ?
\end{enumerate}

\section*{Problème 11 --- pH et échelle logarithmique \hfill [\etoile\etoile\etoile\etoile]}

Deux solutions $S_1$ et $S_2$ ont respectivement des pH égaux à 3,0 et 5,0.

\begin{enumerate}[label=\arabic*)]
    \item Rappeler la définition du pH.
    \item Calculer les concentrations en ions oxonium correspondantes.
    \item Comparer quantitativement les acidités des deux solutions.
    \item Expliquer pourquoi dire « la solution de pH 3 est deux fois plus acide que celle de pH 5 » est faux.
    \item Une dilution par 10 d'une solution acide modifie-t-elle fortement ou faiblement son pH ? justifier.
    \item Pourquoi l'échelle de pH est-elle particulièrement adaptée à la chimie des solutions aqueuses ?
\end{enumerate}

\section*{Problème 12 --- Réactif limitant et vision microscopique \hfill [\etoile\etoile\etoile]}

On mélange $n_0(\ce{Ag^+}) = \SI{3,0e-2}{mol}$ et $n_0(\ce{Cl^-}) = \SI{2,0e-2}{mol}$.
La transformation est modélisée par :
\[
\ce{Ag^+ (aq) + Cl^- (aq) -> AgCl(s)}.
\]

\begin{enumerate}[label=\arabic*)]
    \item Dresser le tableau d'avancement.
    \item Déterminer le réactif limitant.
    \item Déterminer l'avancement maximal.
    \item Déterminer les quantités finales des espèces en solution et du précipité.
    \item Expliquer microscopiquement ce que signifie le fait qu'un réactif soit limitant.
    \item Pourquoi le réactif limitant n'est-il pas nécessairement celui introduit en plus petite quantité de matière dans l'absolu ?
\end{enumerate}

\section*{Problème 13 --- Suivi conductimétrique d'une transformation \hfill [\etoile\etoile\etoile\etoile]}

On suit la réaction entre une solution d'acide chlorhydrique et une solution d'hydroxyde de sodium en mesurant la conductivité du mélange.

\begin{enumerate}[label=\arabic*)]
    \item Pourquoi la conductivité d'une solution dépend-elle de la présence d'ions ?
    \item Expliquer qualitativement l'évolution de la conductivité avant l'équivalence.
    \item Expliquer qualitativement l'évolution de la conductivité après l'équivalence.
    \item Pourquoi l'équivalence peut-elle être repérée à partir d'un changement de comportement de la courbe ?
    \item Comparer l'intérêt d'un suivi conductimétrique et d'un suivi pH-métrique selon le contexte.
\end{enumerate}

\section*{Problème 14 --- Pile Daniell : fonctionnement et interprétation \hfill [\etoile\etoile\etoile\etoile]}

On réalise une pile Daniell à partir des couples $\ce{Zn^{2+}/Zn}$ et $\ce{Cu^{2+}/Cu}$.

\begin{enumerate}[label=\arabic*)]
    \item Écrire les demi-équations électroniques associées aux deux couples.
    \item En déduire l'équation de la réaction globale spontanée.
    \item Identifier l'anode et la cathode.
    \item Indiquer le sens de circulation des électrons dans le circuit extérieur.
    \item Expliquer le rôle du pont salin.
    \item Pourquoi les électrons ne circulent-ils pas dans les solutions ?
    \item Expliquer en quoi la pile réalise une conversion énergétique.
    \item Décrire l'évolution des masses des électrodes au cours du fonctionnement.
\end{enumerate}

\section*{Problème 15 --- Corrosion du fer et protection \hfill [\etoile\etoile\etoile\etoile]}

Une barrière métallique en acier est exposée à l'air humide pendant plusieurs années.

\begin{enumerate}[label=\arabic*)]
    \item Expliquer pourquoi la corrosion du fer est un phénomène d'oxydoréduction.
    \item Citer les conditions favorisant la corrosion.
    \item Pourquoi l'eau salée accélère-t-elle généralement la corrosion ?
    \item Expliquer le principe d'une protection par peinture.
    \item Expliquer le principe d'une galvanisation au zinc.
    \item Pourquoi une couche protectrice rayée n'a-t-elle pas les mêmes conséquences selon qu'il s'agit d'un dépôt de zinc ou d'une simple peinture ?
\end{enumerate}

\section*{Problème 16 --- Identifier une espèce organique par spectroscopie IR \hfill [\etoile\etoile\etoile]}

On étudie une molécule organique de formule brute $\ce{C3H6O2}$. Son spectre IR présente une bande intense vers $\SI{1700}{cm^{-1}}$, et une bande très large dans la zone $2500-3200\ \si{cm^{-1}}$.

\begin{enumerate}[label=\arabic*)]
    \item Que suggère la bande vers $\SI{1700}{cm^{-1}}$ ?
    \item Que suggère la large bande entre $2500$ et $3200\ \si{cm^{-1}}$ ?
    \item Quelle famille chimique peut-on envisager en priorité ?
    \item Proposer une molécule cohérente avec l'ensemble de ces informations.
    \item Expliquer pourquoi la formule brute seule ne suffit pas à identifier la molécule.
    \item Expliquer pourquoi le spectre IR permet d'identifier des groupes caractéristiques mais pas toujours de lever toute ambiguïté.
\end{enumerate}

\section*{Problème 17 --- Synthèse d'un ester et rendement \hfill [\etoile\etoile\etoile\etoile]}

On réalise la synthèse d'un ester à partir d'un alcool et d'un acide carboxylique. Après purification, on obtient une masse de produit inférieure à celle espérée.

\begin{enumerate}[label=\arabic*)]
    \item Écrire l'équation générale d'une estérification.
    \item Pourquoi cette transformation est-elle souvent lente ?
    \item Pourquoi dit-on qu'elle est limitée ?
    \item Définir le rendement de synthèse.
    \item Citer plusieurs raisons expliquant qu'un rendement expérimental soit inférieur à 100\,\%.
    \item Pourquoi le chauffage à reflux est-il souvent utilisé ?
    \item Pourquoi une odeur agréable du produit n'est-elle pas une preuve suffisante de pureté ?
\end{enumerate}

\section*{Problème 18 --- Boisson énergisante et caféine : lecture critique de données \hfill [\etoile\etoile\etoile\etoile\etoile]}

Une boisson énergétique indique sur son étiquette : « 32 mg de caféine pour 100 mL ». La canette contient 250 mL.
Une publicité suggère qu'elle « améliore immédiatement la performance physique et intellectuelle ».

\begin{enumerate}[label=\arabic*)]
    \item Calculer la masse totale de caféine contenue dans une canette.
    \item Expliquer la différence entre une donnée mesurée et une affirmation publicitaire.
    \item Proposer des questions scientifiques permettant d'évaluer expérimentalement l'effet annoncé.
    \item Pourquoi une corrélation éventuelle entre consommation et performance ne suffit-elle pas toujours à établir une causalité ?
    \item Citer des paramètres expérimentaux qu'il faudrait contrôler dans une étude sérieuse.
    \item Expliquer en quoi cet exercice mobilise non seulement des connaissances scientifiques, mais aussi une lecture critique de l'information.
\end{enumerate}

\section*{Problème 19 --- Panneaux photovoltaïques et bilan énergétique \hfill [\etoile\etoile\etoile\etoile]}

On considère une installation photovoltaïque recevant, à midi, une puissance radiative moyenne de $\SI{800}{W.m^{-2}}$ sur une surface de $\SI{12}{m^2}$. Le rendement des panneaux est estimé à 18\,\%.

\begin{enumerate}[label=\arabic*)]
    \item Calculer la puissance radiative reçue par les panneaux.
    \item En déduire la puissance électrique théorique fournie.
    \item Expliquer ce que signifie physiquement un rendement de 18\,\%.
    \item Où « part » le reste de l'énergie reçue ?
    \item Citer des facteurs réels pouvant faire varier la puissance effectivement produite.
    \item Pourquoi la puissance maximale annoncée pour un panneau ne correspond-elle pas à une production constante au cours d'une journée ?
\end{enumerate}

\section*{Problème 20 --- Sujet de synthèse : quel modèle choisir ? \hfill [\etoile\etoile\etoile\etoile\etoile]}

Pour chacune des situations suivantes, on demande de choisir un modèle pertinent, d'en discuter les hypothèses, puis de préciser quelle grandeur physique ou chimique on chercherait à déterminer en priorité.

\begin{enumerate}[label=\alph*)]
    \item Estimer la durée de chute d'une bille d'acier dans l'air.
    \item Déterminer la concentration d'un détartrant acide du commerce.
    \item Expliquer le fonctionnement d'une pile présente dans une télécommande.
    \item Identifier une molécule organique à partir d'un spectre IR et d'une formule brute.
    \item Étudier l'intérêt d'un casque de vélo lors d'un choc.
    \item Expliquer pourquoi deux étoiles très proches peuvent ne pas être distinguées à l'œil nu.
\end{enumerate}

Pour chaque cas, on attend :
\begin{itemize}[leftmargin=1.5em]
    \item le choix d'un modèle ;
    \item les hypothèses ;
    \item les limites du modèle ;
    \item la ou les grandeurs à mesurer ou calculer ;
    \item une courte conclusion scientifique.
\end{itemize}

\hrule

\section*{Conseils de progression}

On peut utiliser les niveaux de difficulté pour organiser le travail :
\begin{itemize}[leftmargin=1.5em]
    \item commencer par les problèmes à \etoile\etoile\etoile{} pour consolider les notions et les modèles classiques ;
    \item poursuivre avec les problèmes à \etoile\etoile\etoile\etoile{} pour travailler l'autonomie, la discussion des hypothèses et les liens entre chapitres ;
    \item terminer par les problèmes à \etoile\etoile\etoile\etoile\etoile{} pour s'entraîner à la synthèse, à la critique scientifique et à la rédaction longue.
\end{itemize}

\section*{Remarques pédagogiques}

Cette fiche est volontairement plus proche de \textbf{gros problèmes de réflexion} que d'une banque de calculs rapides. Elle vise à faire travailler :
\begin{itemize}[leftmargin=1.5em]
    \item l'identification du bon chapitre à mobiliser ;
    \item le choix du bon modèle ;
    \item la lecture critique des hypothèses ;
    \item la capacité à relier plusieurs notions dans un même exercice ;
    \item la rédaction scientifique claire.
\end{itemize}

\end{document}